```latex
\documentclass[11pt]{article}

\usepackage{amsmath,amssymb,amsthm,mathtools}
\usepackage[margin=1in]{geometry}

\newtheorem{theorem}{Theorem}
\newtheorem{lemma}[theorem]{Lemma}
\newtheorem{proposition}[theorem]{Proposition}
\newtheorem{definition}[theorem]{Definition}
\newtheorem{remark}[theorem]{Remark}

\newcommand{\DD}{\mathcal D}
\newcommand{\JJ}{\mathcal J}
\newcommand{\RR}{\mathcal R}
\newcommand{\cN}{\mathcal N}
\newcommand{\om}{\omega}
\newcommand{\eps}{\varepsilon}
\newcommand{\C}{\mathbb C}

\begin{document}

\title{A criterion for holomorphic normal-crossing coordinates near a symplectic SNC divisor}
\author{}
\date{}
\maketitle

\begin{abstract}
The statement in the question is stronger than the standard symplectic
regularization theorem for simple normal-crossing symplectic divisors.  The
standard theorem gives, after a cohomologous deformation of the symplectic
form, compatible tubular neighborhoods and hence almost complex structures for
which local normal coordinates are \(\bar\partial\log\)-compatible:
\[
        \bar\partial_J z_i=z_i\theta_i ,
\]
not necessarily \(J\)-holomorphic.  The difference is exactly the existence of
compatible pseudoholomorphic local frames of the normal line bundles of all
strata.  The theorem below gives a precise necessary and sufficient condition
for the coordinates requested in the question, and proves the usual
regularized construction under that additional normal-holomorphic condition.
\end{abstract}

\section{Set-up}

Let \((M^{2n},\omega)\) be closed and let
\[
        D=\bigcup_{i\in S}D_i
\]
be a symplectic simple normal-crossing divisor.  Thus the \(D_i\) are closed
codimension-two symplectic submanifolds meeting transversely, every stratum
\[
        D_I=\bigcap_{i\in I}D_i,\qquad I\subset S,
\]
is symplectic, and the symplectic orientation of \(D_I\) agrees with the
intersection orientation.  This is the convention of
Tehrani--McLean--Zinger \cite[Definition 2.1]{TMZ}.

The regularization theorem of Tehrani--McLean--Zinger says that, after replacing
\(\omega\) by a cohomologous deformation \(\omega'\) through forms for which
\(D\) remains an SNC symplectic divisor, there exists an \(\omega'\)-regularization
of \(D\).  Concretely, near each stratum \(D_I\) one has compatible
identifications of a neighborhood of \(D_I\) with a neighborhood of the zero
section in
\[
        \bigoplus_{i\in I} N_i \longrightarrow D_I ,
\]
where \(N_i\) is the normal line to \(D_i\) along \(D_I\), equipped with
Hermitian metrics and compatible connections.  See
\cite[Definitions 2.9--2.12 and Theorem 2.13]{TMZ}.

\section{The extra datum needed for genuinely holomorphic normal coordinates}

\begin{definition}[normal-holomorphic regularization]
Let \(\omega'\) be a symplectic form for which \(D\) is an SNC symplectic
divisor, and let \(\RR\) be an \(\omega'\)-regularization.  We say that
\((\omega',\RR)\) is \emph{normal-holomorphic} if the following data are given.

\begin{enumerate}
\item For every stratum \(D_I\), an \(\omega'|_{D_I}\)-tame almost complex
structure \(J_I\).

\item For every \(i\in I\), a complex line-bundle atlas for
\(N_i\to D_I\) whose transition functions are \(J_I\)-holomorphic.

\item These atlases are compatible with the product-Hermitian identifications
between normal bundles supplied by the regularization \(\RR\).
\end{enumerate}
\end{definition}

\begin{theorem}[criterion]
After a cohomologous deformation of \(\omega\) for which \(D\) remains an SNC
symplectic divisor, the almost complex structure requested in the question
exists if the resulting regularization is normal-holomorphic.  Conversely, any
almost complex structure with the local holomorphic-coordinate property of the
question induces a normal-holomorphic structure on the normal line bundles of
all strata.
\end{theorem}

\begin{proof}
We prove the two directions separately.

First assume \((\omega',\RR)\) is normal-holomorphic.  Fix a stratum \(D_I\).
Using the regularization, a neighborhood of \(D_I\) is identified with a
neighborhood of the zero section in
\[
        \bigoplus_{i\in I} N_i \longrightarrow D_I .
\]
Choose a normal-holomorphic trivialization over an open set
\(U\subset D_I\).  Thus the total space is written locally as
\[
        U\times \C^{|I|}
\]
with fiber coordinates \(w_i\), \(i\in I\), and with transition maps
\[
        (x,(w_i))\longmapsto
        \bigl(x,(g_i(x)w_i)\bigr),
\]
where each \(g_i:U\cap U'\to \C^*\) is \(J_I\)-holomorphic.  Define on
\(U\times\C^{|I|}\) the product almost complex structure
\[
        \widehat J_I = J_I\oplus i\oplus\cdots\oplus i .
\]
Because the transition functions \(g_i\) are \(J_I\)-holomorphic, the above
transition maps are \(\widehat J_I\)-holomorphic.  Hence the product almost
complex structures patch to a smooth almost complex structure on a
neighborhood of the zero section in \(\bigoplus_{i\in I}N_i\).

Compatibility of the regularization on overlaps of strata ensures that these
local definitions agree where two such models overlap.  Pushing the resulting
almost complex structure forward by the regularization gives a smooth almost
complex structure \(J\) on a neighborhood \(U_D\) of \(D\subset M\).

At a point lying on precisely the components \(D_i\), \(i\in I\), the local
fiber coordinates \(w_i\) satisfy
\[
        dw_i\circ J=i\,dw_i .
\]
Thus the \(w_i\) are \(J\)-holomorphic and, in these coordinates,
\[
        D\cap U_D = \bigl\{ \prod_{i\in I} w_i=0\bigr\}.
\]

It remains to ensure tameness and extend \(J\) globally.  In the regularized
model, the form \(\omega'\) is equal to the standard model form along the zero
section: it is the direct sum of \(\omega'|_{D_I}\) on the base and the
standard symplectic forms on the complex normal lines.  Since \(J_I\) is
\(\omega'|_{D_I}\)-tame and the fiber complex structures are standard,
\(\omega'(\cdot,J\cdot)>0\) along \(D\).  Tameness is an open condition, so
after shrinking \(U_D\), the constructed \(J\) is \(\omega'\)-tame on \(U_D\).

Choose any global \(\omega'\)-tame almost complex structure \(J_0\) on \(M\).
The bundle over \(M\) whose fiber at \(p\) is the contractible space of
\(\omega'_p\)-tame complex structures on \(T_pM\) has contractible fibers.
Hence the section \(J|_{U_D}\) extends, after shrinking once more and using a
collar interpolation, to a global \(\omega'\)-tame almost complex structure
which equals the constructed one near \(D\).  This gives the desired \(J\).

Conversely, suppose a global \(\omega'\)-tame almost complex structure \(J\)
has the local holomorphic-coordinate property stated in the question.  On a
coordinate chart near a point of a stratum \(D_I\), write the holomorphic
normal coordinates as \(z_i\), \(i\in I\).  The differentials \(dz_i\) identify
the corresponding normal line \(N_i\) with the trivial complex line.  On an
overlap of two such charts the two choices differ by multiplication by a
nowhere-zero function.  Because both coordinate functions are
\(J\)-holomorphic, these transition functions are pseudoholomorphic on the
stratum.  Thus the coordinate charts define pseudoholomorphic local frames of
the normal line bundles.  Restricting \(J\) to each stratum gives the
structures \(J_I\), and the compatibility of the coordinate systems for
normal-crossing charts gives the required compatibility with the normal bundle
identifications.  Therefore a normal-holomorphic regularization is induced.
\end{proof}

\section{What the standard symplectic regularization theorem gives}

The theorem of Tehrani--McLean--Zinger supplies compatible symplectic
regularizations after cohomologous deformation, but it does not by itself
supply normal-holomorphic frames in the above sense.  From a regularization one
canonically obtains almost complex structures adapted to \(D\) for which local
normal coordinates satisfy the weaker logarithmic equation
\[
        \bar\partial_J z_i = z_i\theta_i
\]
for smooth one-forms \(\theta_i\).  This is the
\(\bar\partial\log\)-compatible condition used by Parker
\cite[Definition 14.5]{Parker}.  It is stable under arbitrary smooth
nonvanishing changes of normal coordinates, because if \(z\) satisfies
\(\bar\partial_J z=z\theta\) and \(g\) is smooth and nowhere zero, then
\[
        \bar\partial_J(gz)
        = gz\bigl(\theta+g^{-1}\bar\partial_J g\bigr).
\]
By contrast, genuine \(J\)-holomorphicity of \(z\) is not stable under an
arbitrary smooth change \(z\mapsto gz\); it is stable only when \(g\) is
\(J\)-holomorphic.  This is precisely the extra normal-holomorphic condition
isolated above.

\begin{remark}
Thus the standard deformation-and-regularization theorem proves the existence
of adapted \(\bar\partial\log\)-compatible almost complex structures, not
automatically the stronger local holomorphic-coordinate structure requested in
the question.  To prove the stronger statement in full generality one would
need an additional theorem asserting that the normal line bundles of all strata
admit compatible pseudoholomorphic local frames after the allowed symplectic
deformation.
\end{remark}

\begin{thebibliography}{9}

\bibitem{MS}
D.~McDuff and D.~Salamon,
\emph{Introduction to Symplectic Topology},
3rd ed., Oxford University Press, 2017.

\bibitem{TMZ}
M.~F. Tehrani, M.~McLean, and A.~Zinger,
\emph{Normal crossings singularities for symplectic topology},
Adv. Math. \textbf{339} (2018), 672--748.
See also arXiv:1410.0609.

\bibitem{Parker}
B.~Parker,
\emph{Log geometry and exploded manifolds},
Abh. Math. Semin. Univ. Hambg. \textbf{82} (2012), 43--81.
See especially Section 14.

\end{thebibliography}

\end{document}
```